\section{Related Work}

This section reviews existing literature related to streaming communication in microservice-based systems, with a particular focus on topics that are central to this thesis: streaming communication technologies and interaction models, architectural integration in adaptive microservice middleware, policy enforcement under continuous data exchange, resilience mechanisms for streaming workflows, and validation practices with respect to scalability and performance.

\subsection{Streaming communication technologies and interaction models}

DYNAMOS is implemented as a dynamically composed microservice system deployed on Kubernetes, using gRPC for inter-service communication and RabbitMQ for asynchronous messaging \cite{stutterheim2023thesis,stutterheim2023icsa}. This constrains the design space for introducing streaming functionality to a limited set of technologies that are compatible with containerized microservices and cloud-native deployment models. The primary question addressed in this body of work is which streaming communication mechanisms are suitable for large-scale or continuous data exchange, and under which conditions direct streaming RPC should be preferred over broker-based approaches.

gRPC provides native support for streaming through client-streaming, server-streaming, and bidirectional streaming RPCs. These interaction patterns differ in message flow control, lifetime of connections, buffering behaviour, and error propagation semantics. In the context of DYNAMOS, the choice of streaming pattern is not merely an API concern but has architectural consequences, as it affects memory usage, cancellation behaviour, and the ability to overlap communication and computation. The original DYNAMOS thesis identifies client-streaming as a promising mechanism for transferring large datasets incrementally, reducing peak memory pressure at the cost of increased handling complexity \cite{stutterheim2023thesis}.

However, the use of long-lived streaming connections introduces operational challenges that are largely absent in unary request-response interactions. Starck and Ghofrani show that gRPC streaming complicates load balancing in Kubernetes environments, as persistent connections can cause traffic to become unevenly distributed across replicas \cite{starck2023streaminglb}. This observation is particularly relevant for systems such as DYNAMOS that rely on horizontal scaling and ephemeral execution units. Complementary benchmarking work demonstrates that the performance characteristics of streaming gRPC are sensitive to runtime configuration and virtualization overhead, with measurable effects on latency and resource utilisation in containerized environments \cite{sodankoor2026grpcbench}.

An alternative to direct service-to-service streaming is to externalise the data plane to a messaging or event streaming platform. Broker-based communication decouples producers and consumers in time, delegates buffering and persistence to the broker, and can simplify fan-out scenarios. Comparative benchmarks of Kafka, RabbitMQ, NATS, and ActiveMQ Artemis show that different brokers dominate under different workload characteristics, reinforcing that broker choice is highly workload dependent \cite{ibrahim2026brokers}. A cloud-native literature review comparing Kafka, Pulsar, and RabbitMQ further highlights trade-offs between throughput, reliability, operational complexity, and Kubernetes integration \cite{nuikka2021cloudnative}. While these studies are not conducted in the context of policy-driven middleware, they provide structured criteria that are directly applicable when evaluating broker-based alternatives for DYNAMOS.

Several studies also emphasize that streaming design is influenced not only by technology choice but by the underlying communication model. DataXc contrasts push-based streaming approaches with pull-based designs and reports improved latency stability and throughput under consumer-controlled flow \cite{coviello2022dataxc}. Although the application domain differs, the distinction between producer-driven and consumer-driven streaming is directly relevant to backpressure management and resource utilisation in DYNAMOS-compatible workflows.

Overall, existing literature provides a broad comparison of streaming technologies and interaction patterns, but these comparisons are typically conducted in generic microservice or stream-processing settings. Existing studies are still useful methodologically because they emphasize identical workloads, common measurement points, latency distribution, throughput, resource pressure, and recovery behaviour as recurring benchmark dimensions \cite{karimov2018benchmarking,Henning2021ScalabilityMetrics,systematicstreaming,ciftci2025ipcmechanisms}. The gap addressed in this thesis is therefore not the invention of a new benchmarking methodology, but the application of those dimensions to dynamically composed, policy-driven middleware with ephemeral execution units.

\subsection{Architectural integration of streaming in adaptive microservice middleware}

Beyond the selection of a streaming mechanism, integrating streaming communication into an adaptive microservice architecture raises architectural questions related to modularity, orchestration, and lifecycle management. DYNAMOS is explicitly designed around dynamic microservice composition, where execution chains are generated and deployed per request based on policy evaluation and system state \cite{stutterheim2023thesis,stutterheim2023icsa}. These chains are short-lived by design, which improves isolation and compliance guarantees but complicates the management of long-lived communication channels.

Communication in DYNAMOS is abstracted through a sidecar-based architecture, decoupling communication concerns from core microservice logic. The sidecar pattern is widely used in cloud-native systems to address cross-cutting concerns such as observability, security, and traffic management. Figueira and Coutinho demonstrate how sidecars can be integrated into self-adaptive microservice architectures on Kubernetes, enabling runtime adaptation without violating service modularity \cite{figueira2024selfadaptive}. While their work does not explicitly address streaming data transfer, it illustrates how communication behaviour can be adapted dynamically at runtime.

Several middleware frameworks operationalise sidecars as first-class runtime components. Figueira and Coutinho discuss sidecar-based runtimes such as Dapr, which provide a protocol-agnostic communication abstraction supporting service invocation, pub/sub, and state management through a sidecar process \cite{figueira2024selfadaptive}. Such runtimes demonstrate that heterogeneous communication paradigms, including streaming and messaging, can be unified behind a stable interface. However, this class of runtimes assumes relatively static service topologies and does not address dynamic chain composition driven by policy evaluation, which limits their applicability as a direct solution for DYNAMOS.

Existing architectural work therefore supports the feasibility of integrating streaming into sidecar-based microservice systems but does not address the combination of long-lived streaming connections with ephemeral, dynamically composed execution chains. This limitation is particularly relevant in systems that combine adaptivity with strict policy enforcement requirements.

\subsection{Policy enforcement and compliance under continuous data exchange}

Policy enforcement is a defining characteristic of DYNAMOS, where access rights, execution locations, and data-flow constraints are evaluated prior to execution and enforced through dynamically composed microservice chains \cite{stutterheim2023thesis,stutterheim2023icsa}. In the current architecture, policy enforcement is implicitly scoped to bounded request-response interactions. Introducing streaming communication challenges this assumption, as data transfer may span longer periods during which policies, system state, or execution context may change.

Neumaier et al. propose a policy-aware architecture for decentralized dataset exchange that emphasizes continuous monitoring and runtime validation against machine-readable policies \cite{neumaier2021policyaware}. While their work does not explicitly target streaming RPC, it introduces an important distinction between one-time policy clearance and ongoing compliance monitoring, which becomes critical when data is exchanged incrementally over long-lived connections.

At the communication layer, several studies propose enforcing security and access control through proxies or sidecars. Thiyagarajan et al. present a proxy-based approach for securing gRPC communication using mTLS, JWT authentication, and RBAC \cite{thiyagarajan2025grpcsecurity}. Although the focus is on security rather than data usage policy, the architectural approach is relevant, as it suggests that enforcement can be applied at communication boundaries and potentially at message granularity.

Despite these contributions, there is limited guidance on how policy enforcement should interact with long-lived streaming connections in systems that dynamically compose execution chains. In particular, existing work does not address how policy violations should be detected or handled mid-stream, or how streams should be terminated or adapted when compliance conditions change. These limitations highlight open challenges when introducing streaming into policy-driven middleware such as DYNAMOS.

\subsection{Resilience, backpressure, and fault handling in streaming workflows}

Streaming-based communication introduces operational challenges that are less pronounced in request-response systems. Long-lived data flows must tolerate partial failures, adapt to fluctuating load, and prevent uncontrolled resource growth through effective backpressure mechanisms. These challenges are compounded in DYNAMOS by dynamic microservice composition and ephemeral execution units.

Reactive microservice architectures emphasize asynchronous communication, non-blocking execution, and fault isolation as mechanisms for improving resilience under variable load \cite{rasheedh2021reactive}. Although this work does not focus on streaming RPC, its principles are directly applicable to streaming workflows, where failures and backpressure must be managed continuously rather than per request.

Studies on high-volume data processing in microservices highlight the importance of explicit flow control and decoupling between producers and consumers. Guntupalli and Vamshi show that event-driven architectures combined with messaging systems can mitigate bottlenecks and stabilize performance under sustained data rates \cite{guntupalli2023highvolume}. While these systems typically rely on brokers rather than direct streaming RPC, they reinforce the need for explicit backpressure mechanisms in streaming data paths.

From a systems perspective, Gan and Delimitrou demonstrate that communication overhead and tail latency dominate the performance of large-scale microservice deployments \cite{gan2018architectural}. Their findings underline that resilience and scalability issues often emerge from inter-service communication rather than computation, an effect that is amplified in streaming scenarios.

Existing literature provides well-established principles for resilience and fault tolerance in microservice systems, but these are typically studied either in request-oriented architectures or broker-based streaming platforms. There is limited work examining how these mechanisms interact with direct streaming RPC in dynamically composed microservice chains, leaving questions about how failures and backpressure should be handled in dynamically composed microservice systems.

\subsection{Validation practices, scalability, and performance evaluation}

Empirical evaluation of streaming communication often focuses on comparing streaming and non-streaming interaction styles in terms of throughput, latency, and resource usage. Studies focused specifically on gRPC streaming show that streaming RPCs outperform unary calls for large payloads and continuous data flows by enabling incremental processing and reducing buffering overhead \cite{roohitavaf2019logplayer,sodankoor2026grpcbench}. Complementary protocol-level comparisons between REST and gRPC further establish that binary, multiplexed transports outperform request-response HTTP under load, which provides additional context for choosing gRPC as the baseline communication layer \cite{restgrpccomparison,lahtevenoja2021communication}.

Scalability in containerized streaming systems has been studied through systematic mapping and empirical evaluations. A mapping study on performance analysis techniques for streaming workloads in containerized microservices identifies network contention, scheduling overhead, and coordination costs as primary scalability bottlenecks \cite{systematicstreaming}. Empirical studies of streaming pipelines deployed on Kubernetes further show that streaming approaches scale more gracefully than request-response models when concurrency increases, provided that backpressure and flow control are properly configured \cite{starck2023streaminglb}.

However, most existing evaluations assume static microservice topologies and fixed execution paths. Prior performance studies rarely consider adaptive middleware that dynamically constructs execution chains based on policy evaluation. As a result, there is limited empirical insight into how streaming communication scales when embedded within dynamically adaptive, policy-driven systems such as DYNAMOS. This lack of combined evaluation of streaming, scalability, and adaptivity motivates the experimental focus of this thesis.

